\subsection*{先减少未知信息，再增加计算}
已知 $x+y=7,xy=10$，要求 $x^2+y^2$，并不一定要先求出两个数。目标对交换 $x,y$ 不变，提示我们改用和与积。展开 $(x+y)^2$ 后立即得到答案 $49-20=29$。

这不是偶然的巧招，而是表示选择：若问题不关心两个数谁是谁，就没必要保留这份区分。二次方程的中心、矩阵的迹和行列式也在压缩类似的对称信息。

\subsection*{例子、猜想与证明分工不同}
做小例子可以发现规律，极端例子可以暴露漏洞，图形可以提示不变量；但它们通常还没有覆盖全部情形。猜想形成后，要把观察变成对任意允许输入有效的推理，并记录用了哪些条件。

例如给定和 $s$ 与积 $p$，实数解的存在要求 $s^2-4p\ge0$。即使目标表达式能顺利算出，原问题也可能根本没有实数对象。计算简化之后仍需回头检查可行性。

建立自己的解题记录时，可以写三段：原表示哪里困难，新表示保留了什么，如何把结果返回原问题并检查边界。能解释为什么换表示，就更容易把方法迁移到下一道题。Lean 可以核验具体恒等式，但不会替你决定哪个表示最有启发性。
